% 09-sema-orchestration.tex — Chapter 9: Semantic Analysis: Orchestration
\chapter{Semantic Analysis: Orchestration}
\label{chap:09-sema-orchestration}
\index{semantic analysis!orchestration}
\index{sema\_resolve\_module}

\pnum
The semantic analysis subsystem transforms a raw parse tree into a fully-annotated,
safety-verified AST ready for code generation.  The orchestrator lives in
\srcref{src/sema.h}; the individual passes live in the headers under
\srcref{src/sema/}.

% -------------------------------------------------------------------------
\section{Two-phase analysis}
\label{sec:sema-twophase}
\index{semantic analysis!two phases}
% -------------------------------------------------------------------------

\pnum
Semantic analysis runs in two ordered phases for each function body:

\begin{description}[leftmargin=3cm, style=nextline]
\item[Phase 1 — Resolve.]
  \cname{sema\_resolve\_stmt()} and \cname{sema\_resolve\_expr()} populate
  \cname{->decl} pointers on \cname{EXPR\_IDENTIFIER} and \cname{EXPR\_CALL}
  nodes, expand UFCS calls, instantiate comptime generics, and expand
  destructuring parameters.  The flag \cname{sema\_walk\_phase} is
  \texttt{false} during this phase.
\item[Phase 2 — Walk.]
  \cname{walk\_stmt()} runs type inference (\cname{sema\_infer\_expr}), value
  range analysis (VRA), bounds checking, and bounded-while verification.  The
  flag \cname{sema\_walk\_phase} is \texttt{true} during this phase.  Checks
  that must only run during walk (e.g.\ array bounds) test this flag before
  emitting errors.
\end{description}

\pnum
The separation prevents resolve-phase operations from seeing walk-phase state
(e.g.\ range information that has not yet been computed) and prevents walk-phase
operations from emitting spurious errors on UFCS-desugared calls that are not
yet complete.

% -------------------------------------------------------------------------
\section{sema\_resolve\_module entry point}
\label{sec:sema-entry}
\index{sema\_resolve\_module}
% -------------------------------------------------------------------------

\pnum
\cname{sema\_resolve\_module(decls, module\_path, arena)}
(\srcref{src/sema.h:833}) is called once per compilation with the flat
declaration list produced by \cname{load\_module}.

\begin{algorithm}
\textbf{sema\_resolve\_module}(\textit{decls, module\_path, arena}):
\begin{enumerate}
\item Store \texttt{arena} in \cname{sema\_arena}; build a fresh \cname{RangeTable}.
\item Retrieve the source text and file from the \cname{ModuleNode} record.
\item Clear the global symbol table and populate it via \cname{sema\_build\_scope}.
\item For each \cname{DECL\_FUNCTION} or \cname{DECL\_PROCEDURE} in \textit{decls}:
  \begin{enumerate}[(a)]
  \item Clear the local symbol table.
  \item Insert parameter symbols (including destructuring expansion).
  \item Apply inline parameter constraints and pre-contracts to the range table.
  \item \textbf{Phase 1}: run \cname{sema\_resolve\_stmt} on the entire body.
  \item \textbf{Phase 2}: set \cname{sema\_walk\_phase = true},
    run \cname{walk\_stmt} on the entire body, then clear the flag.
  \item Run \cname{sema\_check\_function\_linearity(d)} to verify ownership.
  \item Clear the local symbol table.
  \end{enumerate}
\item Run \cname{sema\_check\_no\_mutual\_recursion(decls)}.
\end{enumerate}
\end{algorithm}

% -------------------------------------------------------------------------
\section{Symbol table}
\label{sec:sema-symtab}
\index{symbol table}
\index{Symbol struct}
% -------------------------------------------------------------------------

\pnum
The symbol table is implemented in \srcref{src/sema/scope.h} as two separate
open-hash tables with 4096 buckets:

\begin{center}
\begin{tabular}{ll}
\toprule
\textbf{Table} & \textbf{Contents} \\
\midrule
\cname{sema\_globals[4096]} & Top-level declarations: functions, structs, enums, type aliases \\
\cname{sema\_locals[4096]}  & Per-function: parameters, \lain{var} declarations \\
\bottomrule
\end{tabular}
\end{center}

\pnum
Lookup (\cname{sema\_lookup}) checks locals first, then globals, so function
parameters shadow top-level symbols.

\begin{structdoc}{Symbol}{src/sema/scope.h:33}
\cname{name}      & \cname{char *} & Raw identifier string (arena copy). \\
\cname{c\_name}   & \cname{char *} & C-mangled identifier for emit. \\
\cname{type}      & \cname{Type *} & Symbol type (NULL before inference). \\
\cname{decl}      & \cname{Decl *} & Originating declaration (NULL for implicit locals). \\
\cname{is\_global}& \cname{bool}   & True if in \cname{sema\_globals}. \\
\cname{is\_mutable}& \cname{bool}  & True if declared \lain{var}. \\
\cname{next}      & \cname{Symbol *} & Hash-chain next. \\
\end{structdoc}

\pnum
The hash function is DJB2 (\cname{sema\_hash}) modulo 4096.  All \cname{Symbol}
nodes are allocated in \cname{sema\_arena}.  The ``clear'' operations
(\cname{sema\_clear\_globals}, \cname{sema\_clear\_locals}) reset bucket
pointers to \texttt{NULL}; arena memory is not released.

% -------------------------------------------------------------------------
\section{Block scope}
\label{sec:sema-scope}
\index{scope!block scope}
% -------------------------------------------------------------------------

\pnum
Within a function body, block scope is maintained via a linked list of
\cname{ScopeFrame} nodes.  Each frame saves a full snapshot of the 4096-pointer
\cname{sema\_locals} array.

\begin{ccode}
static void sema_push_scope(void) {
    ScopeFrame *frame = arena_push_aligned(sema_arena, ScopeFrame);
    memcpy(frame->saved_locals, sema_locals, sizeof(sema_locals));
    frame->parent = scope_stack;
    scope_stack = frame;
}
static void sema_pop_scope(void) {
    memcpy(sema_locals, scope_stack->saved_locals, sizeof(sema_locals));
    scope_stack = scope_stack->parent;
}
\end{ccode}

\pnum
Symbols declared inside a block (\cname{STMT\_VAR}, \cname{STMT\_FOR} index)
are inserted into \cname{sema\_locals} during the block.  On pop, the previous
snapshot is restored, making the block-local symbols invisible to subsequent code.

\begin{invariant}
Every \cname{sema\_push\_scope()} call is matched by exactly one
\cname{sema\_pop\_scope()} call.  The orchestrator wraps each function body
pass in push/pop pairs; individual statement handlers do the same for their
sub-blocks.
\end{invariant}

% -------------------------------------------------------------------------
\section{Walk-phase statement handler}
\label{sec:sema-walk}
\index{walk\_stmt}
% -------------------------------------------------------------------------

\pnum
\cname{walk\_stmt(s)} (\srcref{src/sema.h:412}) is the walk-phase driver.  For
each statement kind it:
\begin{itemize}
\item Calls \cname{sema\_infer\_expr} on all sub-expressions to set their
  \cname{->type} fields.
\item Updates the \cname{RangeTable} via \cname{range\_set} and
  \cname{sema\_apply\_constraint}.
\item Pushes/pops in-guards for \cname{STMT\_IF} and \cname{STMT\_WHILE}
  conditions containing \lain{in} expressions.
\item Calls \cname{sema\_verify\_bounded\_while} for \lain{while} loops that
  carry a termination measure.
\item Recursively processes sub-statement lists within push/pop scope pairs.
\end{itemize}

\pnum
Loop widening is performed before and after loop bodies via
\cname{sema\_widen\_loop}, which sets the range of every variable assigned inside
the loop to \cname{range\_unknown()}.  This models the conservative fixpoint
used by the VRA.

% -------------------------------------------------------------------------
\section{Mutual recursion detection}
\label{sec:sema-mrec}
\index{mutual recursion detection}
% -------------------------------------------------------------------------

\pnum
After all functions are processed, \cname{sema\_check\_no\_mutual\_recursion(decls)}
(\srcref{src/sema.h:786}) performs DFS over the static call graph.

\pnum
For each non-generic \cname{DECL\_FUNCTION} (\lain{func}, not \lain{proc}), a
depth-first traversal follows \cname{EXPR\_CALL} edges.  When the DFS
revisits a node already on the current stack, a cycle is detected and error
\diagcode{E011} is emitted.

\pnum
The DFS is bounded by \cname{MREC\_MAX\_NODES = 256}.  Functions beyond this
limit are silently skipped; this is acceptable for the current compiler which
does not target programs with more than 256 pure functions in a call cycle.
\lain{proc} declarations are excluded from cycle detection because procedures
are not required to terminate.

\begin{note}
Generic (\cname{comptime}) functions are excluded from cycle detection because
each instantiation is a distinct function body.  Cycles through instantiations
cannot be checked without tracking per-instantiation call graphs.
\end{note}
